\section*{अध्याय \ref{ch:b2:palaeoenvironments} --- जीवाश्म और पुरापर्यावरण}

\begin{solution}{exo:b2:palaeoenvironments:1}
(क) $112/135 = 0.830$ और $112/150 = 0.747$। (ख) तल पर मृत सीपियाँ वर्षों या उससे अधिक में एकत्र
होती हैं: उनमें वे प्रजातियाँ भी हैं जो वहाँ पहले, या कभी-कभार, रहती थीं, या बहाकर लाई गईं, इसलिए
\omterm{def:b2:palaeoenvironments:assemblages}{मृत्यु-समुच्चय} समुदाय की अनेक पीढ़ियों का औसत होता है। (ग) कोमल-शरीर समूह,
विशेषकर कृमि जैसे अपरदभक्षी और मांसाहारी: कंकाल के बिना वे दफ़न से पहले अपघटित हो जाते हैं।
\end{solution}

\begin{solution}{exo:b2:palaeoenvironments:2}
(क) $5 \times 165 = \qty{825}{kyr}$, लगभग \qty{0.8}{Myr}, 0.7 से 1 मिलियन वर्षों के भीतर।
(ख) $183.73 - 182.77 = \qty{0.96}{Myr}$, और $0.96/0.165 = 5.8$: लगभग छह क्षितिज। (ग) एक
वंशक्रम दस लाख वर्ष से बहुत कम में बदलता है; केवल क्षितिज ही क्रमिक रूपों को अलग करते हैं,
जबकि एक क्षेत्र उन्हें एक साथ मिला देता है।
\end{solution}

\begin{solution}{exo:b2:palaeoenvironments:3}
(क) नेरेइटीज़ \omterm{def:b2:palaeoenvironments:ichnofacies}{इक्नोफ़ेसीज़}: किसी गंदलाव-निक्षेप तंत्र का गहरा समुद्र-तल, मग्नतट-भंग से परे, जहाँ
भोजन दुर्लभ है। (ख) स्कोलिथोस \omterm{def:b2:palaeoenvironments:ichnofacies}{इक्नोफ़ेसीज़}: अग्रतट या तटमुख। (ग) क्रूज़ियाना \omterm{def:b2:palaeoenvironments:ichnofacies}{इक्नोफ़ेसीज़}:
अपतट, सामान्य-मौसम और तूफ़ान तरंग-आधारों के बीच। (घ) ज़ूफ़ाइकोस \omterm{def:b2:palaeoenvironments:ichnofacies}{इक्नोफ़ेसीज़}:
मेसोज़ोइक में, निचला मग्नतट या ढलान, तूफ़ान तरंग-आधार के नीचे।
\end{solution}

\begin{solution}{exo:b2:palaeoenvironments:4}
(क) परत A में स्टेनोहेलाइन समूह हैं: सामान्य समुद्री जल। परत B, प्रजातियों में निर्धन पर
व्यष्टियों में समृद्ध, पादप-मलबे के साथ और किसी स्टेनोहेलाइन समूह के बिना, खारे-मीठे जल का संकेत देती है,
उदाहरण के लिए किसी ज्वारनदमुख का। (ख) अनुपस्थित समूह कई हैं, उसी आयु की समुद्री परतों में सामान्य और
अच्छी तरह परिरक्षित: उनकी संयुक्त अनुपस्थिति एक प्रतिरूप है, जबकि एक जीवाश्म संयोग से या
खराब परिरक्षण से अनुपस्थित हो सकता है। (ग) अवसाद (ज्वारीय या नदी-वाहिकाएँ, कीचड़-आवरण, जड़ें), \omterm{def:b2:palaeoenvironments:fossil-kinds}{चिह्न
जीवाश्म} (कुछ ही प्रजातियों वाली एक क्रूज़ियाना \omterm{def:b2:palaeoenvironments:ichnofacies}{इक्नोफ़ेसीज़}), और परत B का पार्श्व में
समुद्री या नदी-निक्षेपों में संक्रमण।
\end{solution}

\begin{solution}{exo:b2:palaeoenvironments:5}
(क) $P = 1/2$ जब $kh/s = \ln 2$: $s = 0.05 \times 0.1/0.693 = \qty{0.0072}{m.yr^{-1}}$,
\qty{7.2}{mm} प्रति वर्ष। (ख) $s$ आधा करने से घातांक दुगुना होता है: $P = (1/2)^2 = 0.25$, आधा।
(ग) $k$ आधा होता है: $P = e^{-2.5} = 0.082$, न कि $e^{-5} = 0.0067$; $e^{2.5} = 12$ गुना अधिक
संभावित। (घ) मोटी, दृढ़ सीपियों का प्रतिनिधित्व अधिक होता है, और दफ़न जितना धीमा उतना अधिक।
\end{solution}

\begin{solution}{exo:b2:palaeoenvironments:6}
(क) ऊपरी सीमा $183.73 + 0.35 = \qty{184.08}{Ma}$ है: चार्ट का \qty{184.2}{Ma} उसके
बाहर है, \qty{0.12}{Myr} से। (ख) $183.73 - 182.77 = \qty{0.96}{Myr}$। (ग) $6/0.96 =
\qty{6.25}{m}$ प्रति मिलियन वर्ष। (घ) जैव-क्षेत्र जीवाश्मों से परिभाषित होता है, जिनकी उपस्थितियाँ और
विलुप्तियाँ विस्तृत क्षेत्रों में लगभग एक साथ होती हैं; उस समय में जमी चट्टान की मोटाई
संचयन की स्थानीय दर पर निर्भर करती है।
\end{solution}

\begin{solution}{exo:b2:palaeoenvironments:7}
(क) $0.49\,e^{-1046/3704} = 0.369$ और $0.49\,e^{-1166/3704} = 0.358$; माध्य $0.364$। (ख) $120
\times (1 - 0.364) = \qty{76.3}{m}$। (ग) $76.3/(1 - 0.49) = \qty{149.6}{m}$। (घ) गुणक
मान लेता है कि पूरी पुनर्स्थापित परत की छिद्रता सतह की $\varphi_0$ है; पर सतह पर
\qty{150}{m} मोटी परत अपने आधार पर अपने शीर्ष से कम छिद्रिल होती है, इसलिए वही कण
\qty{149.6}{m} से कम भरते हैं।
\end{solution}

\begin{solution}{exo:b2:palaeoenvironments:8}
(क) एक इकाई आयतन में $\varphi$ जल और $1 - \varphi$ कण हैं: $\rho = \varphi
\rho_{\text{जल}} + (1 - \varphi)\rho_{\text{कण}}$, जिससे सूत्र मिलता है; $(2710 -
2573)/(2710 - 1000) = 0.080$। (ख) $0.51\,e^{-956/4545} = 0.413$। (ग) सीमेंटन: छिद्रों में
अवक्षेपित कैल्साइट, जिससे कार्बोनेट संपीड़न के अतिरिक्त गुज़रते हैं, केवल दफ़न की तुलना में कहीं अधिक
छिद्रता हटाता है।
\end{solution}

\begin{solution}{exo:b2:palaeoenvironments:9}
(क) गहरे-समुद्री गंदलाव-निक्षेप (नेरेइटीज़ \omterm{def:b2:palaeoenvironments:ichnofacies}{इक्नोफ़ेसीज़}); निचले मग्नतट या ढलान का एक जुरैसिक तल
(ज़ूफ़ाइकोस); तरंग-आधारों के बीच एक अपतट तल (क्रूज़ियाना); एक तटमुख (स्कोलिथोस); भूमि पर एक
दलदल। (ख) जल ऊपर की ओर उथला हुआ: तट समुद्र की ओर खिसका। (ग) समुद्र-स्तर का एक पतन,
या अवधसान से तेज़ अवसाद-आपूर्ति, जैसे किसी बढ़ते डेल्टा के सामने; पतन
उसी समय दूसरी द्रोणियों में भी मिलेगा, डेल्टा स्थानीय होगा और उसका अवसाद एक ओर से
आएगा। (घ) एक अनुरूप अनुक्रम में पड़ोसी पर्यावरण एक-दूसरे के बाद आते हैं; बीच वालों
के ऊपर से छलाँग समय में किसी अंतराल या किसी भ्रंश की ओर संकेत करती है।
\end{solution}

\begin{solution}{exo:b2:palaeoenvironments:10}
(क) दफ़न तक का समय $h/s$ = 100, 10 और \qty{1}{yr} है। कोमल शरीर: $e^{-100} \approx
\num{4e-44}$, $e^{-10} = \num{4.5e-5}$, $e^{-1} = 0.37$। सीपी ($k = \qty{0.01}{yr^{-1}}$):
$e^{-1} = 0.37$, $e^{-0.1} = 0.90$, $e^{-0.01} = 0.99$। (ख) $s = kh/\ln 100 = 0.1/4.605 =
\qty{0.022}{m.yr^{-1}}$, \qty{22}{mm} प्रति वर्ष। (ग) महीनों तक, शायद ही; पर एक तूफ़ानी परत, एक
गंदलाव-निक्षेप या एक राख-पात तल को घंटों में दबा देता है। (घ) अवसादन की स्थिर दर: तब
दफ़न तक का समय माध्य दर वाले $h/s$ से बहुत छोटा होता है, और $P = e^{-kT}$, जहाँ $T$
दफ़न तक का वास्तविक समय है, एक की ओर प्रवृत्त होता है जब $T$, $1/k$ से बहुत कम हो।
\end{solution}

\begin{solution}{exo:b2:palaeoenvironments:11}
(क) एक पतली परत $\mathrm dz$, गहराई $z$ पर, कण $(1 - \varphi(z))\,\mathrm dz$ रखती है;
सतह पर वे $\mathrm dz_0$ भरते हैं, जहाँ $(1 - \varphi_0)\,\mathrm dz_0 = (1 -
\varphi(z))\,\mathrm dz$। (ख) $\varphi = 0.63\,e^{-816/1960} = 0.415$; $(1 - 0.415)/0.37 =
1.581$, इसलिए \qty{158.1}{m}। (ग) \qty{9.5}{m} अधिक: \cref{exo:b2:palaeoenvironments:7} की तरह,
पुनर्स्थापित परत सर्वत्र अपनी सतह जितनी छिद्रिल नहीं है। (घ) कणों की दर,
$h(1-\varphi)/\Delta t$, केवल वर्तमान मोटाई और छिद्रता का उपयोग करती है, दोनों मापी गई।
\end{solution}

\begin{solution}{exo:b2:palaeoenvironments:12}
(क) $9.5/0.165 = 57.6$ क्षितिज। (ख) आज चूना-पत्थर के $180/57.6 = \qty{3.1}{m}$, ताज़े अवसाद के $338/57.6 =
\qty{5.9}{m}$। (ग) अधिक से अधिक \qty{1.5}{m}: $180/1.5 = 120$ अंतराल, 121
नमूने। (घ) निक्षेपण में विराम, ऐसी परतें जहाँ अमोनाइट अनुपस्थित हैं, और कुछ सेंटीमीटर में संघनित
क्षितिज जहाँ अवसाद धीरे आया।
\end{solution}

\begin{solution}{pb:b2:palaeoenvironments:1}
\textbf{1.} $866 \sin \ang{12} = 866 \times 0.2079 = \qty{180.0}{m}$।
\textbf{2.} $481 \times 0.2079 = \qty{100.0}{m}$।
\textbf{3.} बलुआ पत्थर, अध्यारोपण से: भुजाएँ केवल \ang{12} नत हैं, इसलिए परतें
उलटी नहीं हैं और सबसे निचली सबसे पुरानी है।
\textbf{4.} $866 - 766 = \qty{100}{m}$ और $1046 - 866 = \qty{180}{m}$।
\textbf{5.} वहाँ कम से कम \qty{766}{m}; उसका मूल शीर्ष अपरदित है, इसलिए उसकी पूरी मोटाई
अज्ञात है।
\textbf{6.} सतह पर 0.70; $0.70\,e^{-766/1408} = 0.406$, \qty{766}{m} पर।
\textbf{7.} वे स्टेनोहेलाइन समूह हैं: जल में खुले समुद्र की लवणता थी।
\textbf{8.} एक क्रूज़ियाना \omterm{def:b2:palaeoenvironments:ichnofacies}{इक्नोफ़ेसीज़}: अपतट, सामान्य-मौसम और तूफ़ान तरंग-आधारों
के बीच।
\textbf{9.} तूफ़ान: उनकी लहरें तल तक पहुँचीं, सीपियाँ तोड़ीं और उन्हें अपरदनकारी आधारों वाली
परतों में जमा किया।
\textbf{10.} तल का जल ऑक्सीजन में निर्धन था: थोड़े बिल खोदने वाले, पटल बचे रहे, पायराइट बना।
\textbf{11.} ज़ूफ़ाइकोस \omterm{def:b2:palaeoenvironments:ichnofacies}{इक्नोफ़ेसीज़}, मेसोज़ोइक में निचला मग्नतट या ढलान, तूफ़ान तरंग-आधार
के नीचे; कभी-कभी बिल खोदने वालों के लिए पर्याप्त ऑक्सीजन।
\textbf{12.} नहीं: अमोनाइट ऊपर के जल में तैरते थे; मृत्यु के बाद उनके कवच तल पर डूबे या बहकर
पहुँचे और एक \omterm{def:b2:palaeoenvironments:assemblages}{मृत्यु-समुच्चय} बनाते हैं।
\textbf{13.} $184.2 - 174.7 = \qty{9.5}{Myr}$।
\textbf{14.} $180/9.5 = \qty{18.9}{m}$ प्रति मिलियन वर्ष।
\textbf{15.} $(2710 - 2573)/(2710 - 1000) = 0.080$।
\textbf{16.} $180 \times (1 - 0.080)/(1 - 0.51) = \qty{338}{m}$।
\textbf{17.} $338/9.5 = \qty{35.6}{m}$ प्रति मिलियन वर्ष।
\textbf{18.} $180 \times 0.92/9.5 = \qty{17.4}{m}$ प्रति मिलियन वर्ष।
\textbf{19.} शेल ने 11.3 और \qty{4.4}{m} प्रति मिलियन वर्ष दिए: कणों की दर
लगभग चार गुना गिरी। तल के गहरा होने के साथ उथले टीले और मग्नतट का कार्बोनेट आना बंद हो गया,
और केवल कीचड़ की एक धीमी वर्षा उस तक पहुँची।
\textbf{20.} वह आगे बढ़ा: नदियाँ, फिर एक उथला टीला, एक अपतट तल, तूफ़ान तरंग-आधार से नीचे का
एक तल, फिर खड़िया का खुला समुद्र। स्तंभ एक ट्रांसग्रेशन दर्ज करता है।
\textbf{21.} चूना-पत्थर के दौरान $57 - 41 = \qty{16}{m}$, शेल के दौरान $86 - 57 = \qty{29}{m}$,
कुल \qty{45}{m}।
\textbf{22.} $(60 - 5) + 338 + 148.6 = \qty{542}{m}$, $184.2 - 161.5 = \qty{22.7}{Myr}$ में:
\qty{23.9}{m} प्रति मिलियन वर्ष।
\textbf{23.} $45/542 = 0.083$: लगभग \qty{8}{\percent}, और कम यदि चक्र-चार्ट
आयाम को बढ़ा-चढ़ाकर दिखाता है। शेष तल के अवधसान से आया।
\textbf{24.} वार्डन चूना-पत्थर की चूने की कीचड़ जमते समय लगभग
\textbf{\qty[text-series-to-math=true]{36}{m}} प्रति मिलियन वर्ष की दर से संचित हुई।
\end{solution}
